\documentclass[UTF8]{ctexart}
\usepackage{xeCJK}
\usepackage{geometry}
\geometry{left=2cm,right=2cm,top=2.5cm,bottom=2.5cm}
\setlength{\headheight}{12.7pt}

\usepackage{titlesec}
\titleformat{\section}
  {\normalfont\large\bfseries}
  {\thesection}
  {1em}
  {}
\titlespacing*{\section}
  {0pt}
  {3.5ex plus 1ex minus .2ex}
  {2.3ex plus .2ex}

\usepackage{amsmath}
\usepackage{amssymb}
\usepackage{amsthm}
\usepackage{enumitem}
\setlist[enumerate]{itemsep=0pt,topsep=0pt,parsep=0pt,leftmargin=0pt,itemindent=2.5em}
\setlist[itemize]{itemsep=0pt,topsep=0pt,parsep=0pt,leftmargin=0pt,itemindent=2.5em}
\usepackage{fancyhdr}
\pagestyle{fancy}
\fancyhf{}
\usepackage{graphicx}
\usepackage{hyperref}
\usepackage{indentfirst}
\usepackage{lmodern}


\begin{document}
\setlength{\abovedisplayskip}{1pt}
\setlength{\belowdisplayskip}{1pt}

\section{轨道与稳定子}

\hypertarget{sec:定义1.1}{}
\noindent\textbf{定义1.1 (轨道)}

给定\href{01 群作用与 G-集#nameddest=sec:定义1.1}{群作用}
$\cdot:G\curvearrowright X$. 对任意的$x\in X$, 它生成的子$G$-集是
\[
    \langle x\rangle=\{gx\mid g\in G\},
\]
记为$\mathrm{orb}(x)$, 称作$x$在该作用下的\textbf{轨道}.

\vspace{10pt}

\hypertarget{sec:定理1.1}{}
\noindent\textbf{定理1.1}

任给群作用$\cdot:G\curvearrowright X$, 所有轨道的全体
\[
    \,_{\displaystyle G}\backslash\!^{\displaystyle X}
    \overset{\mathrm{def}}{=}\{\mathrm{orb}(x)\mid x\in X\}
\]
组成$X$的一个划分.

\noindent{\kaishu 证明.}

\begin{enumerate}
    \item 显然任意的$\mathrm{orb}(x)$非空,
    \item 如果$\mathrm{orb}(x_1)=\mathrm{orb}(x_2)$,
    再假设存在$x\in\mathrm{orb}(x_1)\cap\mathrm{orb}(x_2)$,
    那么存在$g_1,g_2\in G$使得
    \[
        x=g_1x_1=g_2x_2\quad\implies\quad x_1=g_1^{-1}g_2x_2\in\mathrm{orb}(x_2),
    \]
    于是$\mathrm{orb}(x_1)\subset\mathrm{orb}(x_2)$,
    同理$\mathrm{orb}(x_2)\subset\mathrm{orb}(x_1)$, 矛盾,
    \item 对任意的$x\in X$, 有$x\in\mathrm{orb}(x)$. \hfill $\qedsymbol$
\end{enumerate}

\vspace{10pt}

\hypertarget{sec:定义1.2}{}
\noindent\textbf{定义1.2}

给定群作用$\cdot:G\curvearrowright X$和$x\in X$, 定义
\[
    \mathrm{Stab}(x)=\{g\in G\mid gx=x\},
\]
称之为$x$在该作用中的\textbf{稳定子}, 它是$G$的子群.

\noindent{\kaishu 验证.}

任取$g_1,g_2\in\mathrm{Stab}(x)$, 有$g_1g_2^{-1}x=g_1g_2^{-1}g_2x=g_1x=x$,
即$g_1g_2^{-1}\in\mathrm{Stab}(x)$. \hfill $\qedsymbol$

\vspace{10pt}

\hypertarget{sec:定理1.2}{}
\noindent\textbf{定理1.2}

任给群作用$\cdot:G\curvearrowright X$, $g\in X$和$x\in X$, 有
\[
    \mathrm{Stab}(gx)=g\,\mathrm{Stab}(x)\,g^{-1}.
\]

\noindent{\kaishu 证明.}

一方面, 对任意的$g'\in\mathrm{Stab}(gx)$, 有$g^{-1}g'gx=g^{-1}gx=x$,
从而$g'=g(g^{-1}g'g)g^{-1}\in g\,\mathrm{Stab}(x)\,g^{-1}$.

另一方面, 对任意的$h\in\mathrm{Stab}(x)$, 有$ghg^{-1}gx=ghx=gx$,
即$ghg^{-1}\in\mathrm{Stab}(gx)$. \hfill $\qedsymbol$

\vspace{10pt}

\hypertarget{sec:定义1.3}{}
\noindent\textbf{定义1.3}

给定群作用$\cdot:G\curvearrowright X$, 如果所有$x\in X$同属一个轨道, 即
\[
    \forall x,x'\in X\,\exists g\in G:x'=gx,
\]
就称$\cdot$是\textbf{传递}的. 如果
\[
    \forall x\in X:\quad\mathrm{Stab}(x)=\{1_G\},\\[-3pt]
\]
即只有$1_G$保持某元素不动, 就称$\cdot$是\textbf{自由}的.
显然这时它也是忠实的.





\newpage

\hypertarget{sec:定理1.3}{}
\noindent\textbf{定理1.3 (轨道-稳定子定理)}

任给群作用$\cdot:G\curvearrowright X$和$x\in X$,
对任意的$y\in\mathrm{orb}(x)$定义
\[
    \varphi(y)=\{g\in G\mid gx=y\},\\[3pt]
\]
则$\varphi:\mathrm{orb}(x)\cong\,\!^{\displaystyle G}\!
/_{\displaystyle\mathrm{Stab}(x)}$, 是这两个$G$-集之间的等变同构.

\vspace{5pt}
特别地, $|\mathrm{orb}(x)|=[G:\mathrm{Stab}(x)]$.

\vspace{3pt}
\noindent{\kaishu 证明.}

首先, 对任意的$g,h\in G$, 有
\[
    h\in\varphi(gx)\iff g^{-1}hx=g^{-1}gx=x\iff
    g^{-1}h\in\mathrm{Stab}(x)\iff h\in g\,\mathrm{Stab}(x),\\[3pt]
\]
即$\varphi(gx)=g\,\mathrm{Stab}(x)$.
于是有$\varphi:\mathrm{orb}(x)\to\,\!^{\displaystyle G}\!
/_{\displaystyle\mathrm{Stab}(x)}$且是满射.

\vspace{3pt}
另外, 如果$\varphi(y)=\varphi(y')$, 取其中的一个$h$,
则$y=hx=y'$. 所以$\varphi$是双射.

最后, 对任意的$g\in G$和$y\in\mathrm{orb}(x)$, 取一个$h\in\varphi(y)$, 则有
\[
    \varphi(gy)=\varphi(ghx)=gh\,\mathrm{Stab}(x)=g\varphi(hx)=g\varphi(y),
\]
从而$\varphi$是等变映射. \hfill $\qedsymbol$

\vspace{15pt}

由轨道-稳定子定理, 同一轨道中不同的$x$给出的陪集空间$\,\!^{\displaystyle G}\!
/_{\displaystyle\mathrm{Stab}(x)}$之间是同构的.

\vspace{3pt}
如果$\cdot:G\curvearrowright X$是传递的, 那么任取$x\in X$有
$X=\mathrm{orb}(x)\cong\,\!^{\displaystyle G}\!/_{\displaystyle\mathrm{Stab}(x)}$.

\vspace{15pt}

\hypertarget{sec:定理1.4}{}
\noindent\textbf{定理1.4}

任给群作用$\cdot:G\curvearrowright X$,
对每个轨道$o\in\,_{\displaystyle G}\backslash\!^{\displaystyle X}$
选取$x_o\in o$. 如果$G,X$都有限, 那么
\vspace{5pt}
\[
    \frac{|X|}{|G|}=\sum_{o\in G\backslash X}\frac{1}{|\mathrm{Stab}(x_o)|}.
\]

\noindent{\kaishu 证明.}

$\displaystyle|X|=\sum_{o\in G\backslash X}|o|=
\sum_{o\in G\backslash X}\frac{|G|}{|\mathrm{Stab}(x_o)|}$. \hfill $\qedsymbol$

\vspace{15pt}

\hypertarget{sec:定理1.5}{}
\noindent\textbf{定理1.5 (Burnside引理)}

任给群作用$\cdot:G\curvearrowright X$, 如果$G,X$都有限, 那么轨道数
\vspace{3pt}
\[
    \left|\,_{\displaystyle G}\backslash\!^{\displaystyle X}\right|=
    \frac{1}{|G|}\sum_{g\in G}|X^g|,\\[-3pt]
\]
其中$X^g\overset{\mathrm{def}}{=}\{x\in X\mid gx=x\}$, 称作$g$的不动点集.

\noindent{\kaishu 证明.}

考虑$T=\{(g,x)\in G\times X\mid gx=x\}$.
根据\href{../../../01 基础集合论/03 自然数/04 有限求和/03 有限 Fubini 定理#nameddest=sec:定理1.1}
{有限Fubini定理}, 一方面,
\[
    |T|=\sum_{g\in G}|X^g|,
\]
另一方面,
\[
    |T|=\sum_{x\in X}|\mathrm{Stab}(x)|=
    \sum_{o\in G\backslash X}\sum_{x\in o}\frac{|G|}{|\mathrm{orb}(x)|}=
    |G|\sum_{o\in G\backslash X}|o|\cdot\frac{1}{|o|}=
    |G|\cdot\left|\,_{\displaystyle G}\backslash\!^{\displaystyle X}\right|,
\]
所以$\displaystyle\left|\,_{\displaystyle G}\backslash\!^{\displaystyle X}\right|=
\frac{1}{|G|}\sum_{g\in G}|X^g|$. \hfill $\qedsymbol$

\end{document}